% !TeX root = ../main.tex
% Lezione 4.pdf, pp. 2--5. Le precisazioni sulle sezioni e sui quantificatori
% completano la formulazione sintetica del manoscritto.
\section{I limiti dei razionali: l'irrazionalità della radice di due}
% Fonte: Lezione 4, p.~2

I razionali permettono di eseguire le quattro operazioni, con la sola
esclusione della divisione per zero, ma non risolvono ogni equazione
algebrica. L'esempio fondamentale è $x^2=2$: la sua soluzione positiva,
indicata con $\sqrt2$, non è razionale.

\begin{proposizione}
Non esiste $d\in\Q$ tale che $d^2=2$; in particolare $\sqrt2\notin\Q$.
\end{proposizione}
\begin{proof}
Supponiamo per assurdo che $d=p/q$, con $p,q\in\Z$, $q\ne0$ e
la frazione ridotta ai minimi termini, soddisfi $d^2=2$. Si avrebbe
\[
  p^2=2q^2.
\]
Quindi $p^2$ è pari e anche $p$ deve essere pari: infatti il quadrato
di un dispari $2h+1$ è $4h(h+1)+1$, che è dispari. Scriviamo allora
$p=2k$. Sostituendo nella relazione precedente otteniamo
\[
  4k^2=2q^2,\qquad q^2=2k^2.
\]
Anche $q^2$, e dunque $q$, è pari. Ma allora $p$ e $q$ sono entrambi
divisibili per $2$, in contraddizione con la scelta di una frazione
ridotta ai minimi termini.
\end{proof}

\section{I reali e l'assioma di Dedekind}
% Fonte: Lezione 4, pp.~3--4

Per descrivere la retta numerica si introduce il campo dei numeri reali
$\R$, che contiene $\Q$. La somma e il prodotto sono operazioni interne
e l'ordine $\le$ è totale. Le proprietà algebriche e di compatibilità
con l'ordine sono le seguenti:
\begin{enumerate}
  \item la somma è associativa e commutativa e ha elemento neutro $0$;
  \item ogni $x\in\R$ ha un opposto $-x$, con $x+(-x)=0$;
  \item il prodotto è associativo e commutativo e ha elemento neutro $1$,
    con $1\ne0$;
  \item ogni $x\ne0$ ha un reciproco $x^{-1}$, con $xx^{-1}=1$;
  \item il prodotto è distributivo rispetto alla somma:
    $x(y+z)=xy+xz$;
  \item se $a\le b$, allora $a+c\le b+c$ per ogni $c\in\R$;
  \item se $a\le b$ e $c\ge0$, allora $ac\le bc$.
\end{enumerate}
Queste proprietà fanno di $\R$ un campo ordinato. La proprietà che lo
distingue da $\Q$ è la \emph{completezza}, espressa qui dall'assioma
di Dedekind.

\begin{assioma}[Dedekind]
\label{ass:dedekind}
Siano $A,B\subseteq\R$ non vuoti e tali che $a\le b$ per ogni
$a\in A$ e ogni $b\in B$. Allora esiste $c\in\R$ per cui
\[
  a\le c\le b\qquad\text{per ogni }a\in A\text{ e ogni }b\in B.
\]
\end{assioma}

Il numero $c$ separa i due insiemi: si trova a destra di tutti gli
elementi di $A$ e a sinistra di tutti quelli di $B$. L'ordine dei
quantificatori richiede un valore di $c$ che funzioni
contemporaneamente per tutte le coppie:
\[
  \exists c\in\R\ \forall a\in A\ \forall b\in B:\ a\le c\le b.
\]
L'assioma afferma l'esistenza, senza richiedere che il separatore
sia determinato in modo univoco. Per esempio,
se $A=\{0\}$ e $B=\{1\}$, ogni $c\in[0,1]$ è un separatore.

\subsection{Perché l'assioma fallisce nei razionali}

Consideriamo gli insiemi
\[
  A=\{x\in\Q:x\ge0,\ x^2\le2\},\qquad
  B=\{x\in\Q:x\ge0,\ x^2\ge2\}.
\]
Sono entrambi non vuoti: $1\in A$ e $2\in B$. Inoltre $a\le b$
per ogni $a\in A$ e $b\in B$, perché sono non negativi e
$a^2\le2\le b^2$. Tuttavia nessun razionale può separare $A$ e $B$.

Per verificarlo, un eventuale separatore razionale $c$ dovrebbe
soddisfare $1\le c\le2$. Se $c^2<2$, scegliamo un razionale
\[
  0<h<\min\left\{1,\frac{2-c^2}{2c+1}\right\}.
\]
Allora $(c+h)^2<c^2+(2c+1)h<2$, quindi $c+h\in A$ ma $c+h>c$:
il separatore non sarebbe a destra di tutti gli elementi di $A$.
Se invece $c^2>2$, scegliamo un razionale
\[
  0<h<\min\left\{1,\frac{c^2-2}{2c}\right\}.
\]
Si ha $c-h>0$ e $(c-h)^2>c^2-2ch>2$, perciò $c-h\in B$ ma
$c-h<c$, ancora una contraddizione. Rimane soltanto la possibilità
$c^2=2$, esclusa dall'irrazionalità appena dimostrata. In $\R$ il
separatore è il punto $\sqrt2$ che manca a $\Q$.

\subsection{L'idea delle sezioni di Dedekind}
% Fonte: Lezione 4, p.~5

È possibile costruire i reali a partire dai razionali mediante le
\emph{sezioni di Dedekind}. Si considera una coppia $(A,B)$ di sottoinsiemi non vuoti
di $\Q$ che soddisfi
\[
  A\cap B=\varnothing,\qquad A\cup B=\Q,\qquad
  a\le b\quad\forall a\in A,\ \forall b\in B.
\]
La coppia descrive un \emph{taglio} della retta razionale in una parte
inferiore e una superiore. Poiché le due parti sono disgiunte,
in realtà $a<b$ per ogni $a\in A$, $b\in B$.

Per rendere la rappresentazione univoca occorre fissare una convenzione
sull'estremo: scegliamo che $A$ non abbia massimo. Infatti, per un
razionale $r$, sia la coppia
$\bigl(\{q:q<r\},\{q:q\ge r\}\bigr)$ sia la coppia
$\bigl(\{q:q\le r\},\{q:q>r\}\bigr)$ soddisfano le tre condizioni
precedenti, pur descrivendo lo stesso punto. La convenzione sceglie
la prima. Con questa precisazione, tali sezioni possono essere usate,
con opportune operazioni e ordine, per costruire un modello di $\R$.
Un taglio può rappresentare anche un punto irrazionale: in quel caso
nessuna delle due parti contiene un estremo razionale nel punto
di separazione.
